% Part: computability
% Chapter: recursive-functions
% Section: composition

\documentclass[../../../include/open-logic-section]{subfiles}

\begin{document}

\olfileid{cmp}{rec}{com}
\olsection{ترکیب}

اگر $f$ و~$g$ دو تابع یک‌موضعی روی اعداد طبیعی باشند، می‌توانیم آن‌ها
را ترکیب کنیم: $h(x)=f(g(x))$. در این صورت تابع تازهٔ~$h(x)$ با
\emph{ترکیب} از توابع $f$ و~$g$ تعریف شده است. می‌خواهیم این کار را
به توابعِ دارای بیش از یک آرگومان تعمیم دهیم.

یک راه چنین است: فرض کنید $f$ تابعی $k$موضعی و
$g_0$، \dots،~$g_{k-1}$ تعداد $k$ تابع باشند که همگی $n$موضعی‌اند.
آنگاه می‌توان تابع تازه و $n$موضعیِ~$h$ را چنین تعریف کرد:
\[
h(x_0, \dots, x_{n-1}) =
f(g_0(x_0, \dots, x_{n-1}), \dots, g_{k-1}(x_0, \dots, x_{n-1}))
\]
اگر $f$ و همهٔ $g_i$ها محاسبه‌پذیر باشند، $h$ نیز محاسبه‌پذیر است.
برای محاسبهٔ $h(x_0,\dots,x_{n-1})$، نخست به ازای هر
$i=0$، \dots،~$k-1$ مقدار
$y_i=g_i(x_0,\dots,x_{n-1})$ را محاسبه می‌کنیم. سپس این مقدارها را
به~$f$ می‌دهیم تا
$h(x_0,\dots,x_{n-1})=f(y_0,\dots,y_{k-1})$ محاسبه شود.

این شاید توصیفی بیش‌ازاندازه محدود از اتفاقی به نظر رسد که هنگام
محاسبهٔ تابعی تازه با توابع موجود رخ می‌دهد. نخست اینکه گاه همهٔ
آرگومان‌های تابع را به کار نمی‌بریم؛ چنان‌که برای استفاده در تعریف
بازگشتی اولیهٔ~$\Add$، تابع $g(x,y,z)=\Succ(z)$ را تعریف کردیم. فرض
کنید مجاز باشیم از توابع زیر استفاده کنیم:
\[
\Proj{n}{i}(x_0, \dots, x_{n-1}) = x_i
\]
توابع~$\Proj{n}{i}$ را توابع \emph{افکنش} می‌نامند؛
$\Proj{n}{i}$ تابعی $n$موضعی است. آنگاه می‌توان $g$ را چنین تعریف
کرد:
\[
g(x, y, z) = \Succ(\Proj{3}{2}(x, y, z)).
\]
در اینجا تابع یک‌موضعیِ~$\Succ$ نقش $f$ را دارد، پس $k=1$. تابع
سه‌موضعیِ~$\Proj{3}{2}$ نیز نقش $g_0$ را دارد. حاصل تابعی سه‌موضعی
است که جانشین آرگومان سوم را بازمی‌گرداند.

توابع افکنش همچنین امکان می‌دهند با جابه‌جاکردن یا یکی‌گرفتن
آرگومان‌ها توابع تازه‌ای تعریف کنیم. برای نمونه، تابع
$h(x)=\Add(x,x)$ را می‌توان چنین تعریف کرد:
\[
h(x_0) = \Add(\Proj{1}{0}(x_0),\Proj{1}{0}(x_0)).
\]
اینجا $k=2$ و $n=1$؛ $\Add$ نقش $f(y_0,y_1)$ را دارد و نقش‌های
$g_0(x_0)$ و $g_1(x_0)$ را هر دو $\Proj{1}{0}(x_0)$، یعنی تابع
افکنش یک‌موضعی (همان تابع همانی)، بر عهده دارند.

اگر $f(y_0,y_1)$ تابعی باشد که از پیش داریم، می‌توان تابع
$h(x_0,x_1)=f(x_1,x_0)$ را چنین تعریف کرد:
\[
h(x_0, x_1) = f(\Proj{2}{1}(x_0, x_1),\Proj{2}{0}(x_0, x_1)).
\]
اینجا $k=2$ و $n=2$ است و $\Proj{2}{1}$ و~$\Proj{2}{0}$ به‌ترتیب
نقش $g_0$ و~$g_1$ را دارند.

شاید این نیز مایهٔ نگرانی باشد که همهٔ
$g_0$، \dots،~$g_{k-1}$ باید یک آریتیِ~$n$ داشته باشند. (به یاد
آورید که \emph{آریتیِ} تابع شمار آرگومان‌های آن است؛ تابع
$n$موضعی آریتیِ~$n$ دارد.) اما افزودن توابع افکنش انعطاف مطلوب را
فراهم می‌کند. برای نمونه، فرض کنید $f$ و~$g$ توابعی سه‌موضعی و~$h$
تابع دوموضعیِ زیر باشد:
\[
h(x,y) = f(x,g(x,x,y),y).
\]
تعریف~$h$ را می‌توان با توابع افکنش چنین بازنوشت:
\[
h(x,y) = f(\Proj{2}{0}(x,y), g(\Proj{2}{0}(x,y), \Proj{2}{0}(x,y),
\Proj{2}{1}(x,y)), \Proj{2}{1}(x,y)).
\]
پس $h$ ترکیب $f$ با $\Proj{2}{0}$، $l$ و~$\Proj{2}{1}$ است، که در آن
\[
l(x,y) = g(\Proj{2}{0}(x,y),\Proj{2}{0}(x,y),\Proj{2}{1}(x,y)),
\]
یعنی $l$ ترکیب $g$ با $\Proj{2}{0}$، $\Proj{2}{0}$ و~$\Proj{2}{1}$
است.

\end{document}
